% =====================================================================
% Cappello di sezione + sec:c-e (integrata in paper.tex via \input)
%
% Labels definiti qui e usati nelle sottosezioni
% (ordine = sequenza degli \input, NON numeri fissi):
%   sec:c-e        Lo stream minimo            (identity)
%   sec:pointer    La posizione di lettura     (pointer)
%   sec:griglia    La griglia temporale        (distribution)
%   sec:deviazione Deviazione: ampiezza e probabilità (deviation, probability)
%   —              L'esempio di mezzo          (duration)
%   —              L'esempio completo          (complete_example)
%   sec:voci       La molteplicità delle voci  (voices)
%
% TODO prima della submission:
%   [x] link audio anonimo (OSF view-only) nella footnote del cappello
%   [x] i path examples/ esistono dopo la rinomina: verificati sul .fls
%       della compilazione. Nomi = token semantico della cartella, senza
%       numeri d'ordine (convenzione in paper/examples/README.md)
% =====================================================================

% --- note di questa sezione (definite qui, richiamate nel corpo) ---
\newcommand{\notaRepo}{\footnote{La libreria Python (compatibilità da 3.9 a 3.12)
è manutenuta presso un repository
github METTERE LINK.
Dipendenze principali NumPy e \texttt{soundfile} per rendering audio in tempo differito,
matplotlib per la \textsc{map}. Il costo di calcolo non cresce con la durata d'ascolto ma col numero di
grani.}}

\newcommand{\notaEsempi}{\footnote{I listati YAML e le realizzazioni audio di
tutti gli esempi di questo capitolo sono disponibili in un archivio ad accesso
aperto: \emph{https://osf.io/xqmh5/?view\_only=5a6d9ac26adb466697477801e73b153f}.}}

\newcommand{\notaGit}{\footnote{\emph{git} è un sistema di controllo di versione
distribuito: registra la storia completa delle modifiche a un insieme di file
testuali, permette di confrontare due stati (\emph{diff}) e di tornare a
qualsiasi punto precedente. Software libero, documentazione e distribuzione
su \emph{https://git-scm.com}.}}

\newcommand{\notaBande}{\footnote{Finestrare un
grano ne modula l'ampiezza, aggiungendo allo spettro bande
laterali~\cite{Roads2001}. Nella sovrapposizione a fattore 2 questi contributi
si elidono, e la differenza RMS gain-matched fra sorgente e resa non supera i
$-74$~dB, sotto la soglia udibile.}}

\newcommand{\notaGamma}{\footnote{Come GAMMA (prima esecuzione assoluta il 10 luglio 2025
al Museo degli Strumenti Musicali di Roma, nel concerto \emph{INCONTRI} del Festival ArteScienza 2025), 
brano acusmatico con un altro motore di composizione assistita per sintesi additiva.}}

\newcommand{\notaClaude}{\footnote{DA SCRIVERE.}}

% ================================================================
\section{DIRAC: un'architettura per l'indagine parametrica}\label{sec:architettura}
PythonGranularEngine\notaRepo\ è un motore
di sintesi per comporre attraverso granulazione.
È scritto con un fine pedagogico per chiunque voglia acquisire consapevolezza
sulla tecnica come strumento di pensiero compositivo. Si discosta da una prassi
consolidata, che Roads~\cite[p.~32]{Roads2021} descrive come una composizione per brevi
improvvisazioni in ambienti real-time, dove si lavora mentre
«the tape is running» e le parti riuscite si selezionano in un
secondo momento. L'interfaccia in questi ambienti
è spesso un controller~\cite[pp.~193--194]{Roads2001} per colmare il gap
del rapporto con lo strumento e con la sua gestualità e il materiale registrato viene poi 
montato spesso in ambienti esterni a quello di sintesi.

Questa libreria è invece pensata per avere come input un linguaggio di dominio testuale in YAML, che espone
un'interfaccia dichiarativa. La scelta è il risultato di un processo di ricerca personale
di circa due anni attraverso altre esperienze\notaGamma.
I motivi maggiori sono due. Il primo è il versionamento — \emph{git}\notaGit\ come documentazione del processo,
con la riproducibilità e la sostenibilità del processo compositivo che ne seguono — insieme alla possibilità di lavorare
in dialogo con un modello linguistico\notaClaude, che può leggere e ragionare facilmente sull'input testuale e divenire una potenziale
interfaccia alternativa. Il secondo è la possibilità di esporre questo motore come API, per estenderlo
e utilizzarlo con un framework compositivo di livello ancor più alto.

Il listato YAML è una struttura a dizionari: coppie di chiave e valore,
annidabili l'una dentro l'altra.
Porterò una serie di esempi\notaEsempi\ che parte dalla semplice riproduzione
del materiale audio scelto e procede modificando poche chiavi alla volta.
Di ogni esempio do il listato — del quale riporto solo le chiavi che si
scostano dallo stream nella sua dichiarazione minima (\S\ref{sec:c-e}) — e la sua \textsc{map}
(\emph{Multiparametric Audio Plot}), una mappa sinottica
come in Fig.~\ref{fig:c-e}: in ascissa il tempo dello stream, in ordinata la
posizione di lettura nel file, e a sinistra, sullo stesso asse verticale, la
forma d'onda della sorgente.
Ogni grano è un poligono: la posizione
indica orizzontalmente l'\emph{on-set} e verticalmente il punto di lettura,
la larghezza e l'altezza sono rispettivamente la sua durata nel tempo e lo spazio percorso sul buffer.
L'orientamento della testa indica il verso di lettura del buffer, e sulla
testa è disegnato il profilo della
finestra applicata a quel grano: entrambi leggibili nella lente di
Fig.~\ref{fig:c-e}, che ingrandisce tre grani contigui.
Sotto la \textsc{map}, quando ci sono, compaiono gli inviluppi dichiarati, coi
breakpoint annotati ciascuno col proprio valore.
A destra, quando necessario, è plottabile una legenda della scala cromatica
(\emph{colorbar}) che mostra la trasposizione
del singolo grano, che
apparirà così più o meno colorato, come visibile in Fig.~\ref{fig:dimensioni}.
La scala è lineare e l'unità è stampata sulla legenda stessa: cents dove il
range si autoregola sull'escursione dei grani presenti, rapporto di frequenza
dove è invece fissato a priori.


La grafia unita alla specifica YAML dichiarata in input permette di osservare. 
Per economia di mezzi e spazio 
userò quasi sempre un solo stream per volta per non stressare la grandezza della figura.
In esempi più fitti con molti parametri e streams attivi è divenuto
necessario nel mio workflow nello stato di test o debugging e 
per comprendere come agire puntualmente su una 
specifica popolazione di grani.
